
nombre d'éléments d'un système minimum de générateurs de l'idéal correspondant; d'après la démonstration qui précède, on a
\[
 \rang(\underline{N}_{X/Z})=\rang(\underline{N}_{X/Y})+\rang(\underline{N}_{Y/Z}).
\]
Il en résulte que le noyau de $j^*(\underline{N}_{X/Y})\longrightarrow \underline{N}_{X/Z}$ est localement libre de rang nul, donc que la suite (1.6.1) est exacte.

ii) On peut supposer $J/J^2$ libre sur $C$, et $I/I^2$ libre sur $B$; de plus, l'hypothèse faite sur la suite des faisceaux conormaux permet de supposer que $I/I^2\otimes_A A/J$ est facteur direct de $J/J^2$. On peut donc trouver des éléments $f_1,\ldots,f_n$ de $A$ relevant une base de $I/I^2$, des éléments $g_1,\ldots,g_m$ de $A$ tels que $f_1,\ldots,f_n,g_1,\ldots,g_m$ relèvent une base de $J/J^2$, les images de $g_1,\ldots,g_m$ engendrant $K/K^2$; quitte à localiser encore, on peut supposer que $f_1,\ldots,f_n$ forment un système minimum de générateurs de $I$, que $f_1,\ldots,f_n,g_1,\ldots,g_m$ un système minimum de générateurs de $J$, et que $\bar g_1,\ldots,\bar g_m$ engendrent $K$. On peut alors construire le diagramme (1.7.1), et utiliser la suite spectrale (1.7.3).

D'après (1.7.2), $i$ étant régulière, on a $E^2_{pq}=0$ pour $q\geq1$; on obtient donc des edge-isomorphismes : $E^2_{p,0}\simeq\Tor_p^{A_0}(C_0,A)$. Comme $i\circ j$ est régulière, ceci entraine d'après (1.7.2) que $E^2_{p,0}=0$ pour $p\geq1$, et par suite que $j$ est régulière, toujours d'après (1.7.2).

iii) Cette dernière propriété est démontrée dans EGA IV 19.1.5 (iv) puisque dans le cas noethérien la définition de EGA IV coïncide avec celle que nous adoptons ici.

\textbf{Proposition 1.8.} \uline{Soient} \(X\) \uline{un schéma}, \(\underline{E}\) \uline{un} \(\underline{O}_X\)-\uline{Module localement libre de type fini}, \(u:\underline{E}\to\underline{O}_X\) \uline{un homomorphisme régulier}, \(\underline{J}=u(\underline{E})\). \uline{On note} \(Y\)
